\chapter{DeepSeek Harness}
\label{ch:deepseek-harness}
\chapterlead{DeepSeek Harness 是 DeepSeek 官方开源的 Agent 运行时，采用插件化架构，提供终端与 Web 两种界面。本章说明安装、配置与第一次任务，并标注它当前处于开发者预览阶段。}

\section{版本与前提}

截至 2026 年 11 月 25 日，可用的包与要求如下。

\begin{booktable}{\textwidth}{J{0.741} J{1.259}}
\toprule
项目 & 值 \\
\midrule
\endhead
npm 包 & \texttt{@deepseek-ai/dsh} \\
版本 & 0.2.2-rc.5 \\
运行时 & Node.js 24.9 或更新 \\
命令 & \texttt{dsh} \\
许可 & MIT \\
官方文档 & \texttt{deepseek-harness.\allowbreak{}github.io} \\
\bottomrule
\end{booktable}

项目处于开发者预览阶段，官方明确说明会有破坏性变更。命令与选项可能随版本改动，执行前用 \texttt{-{}-help} 核对。

\begin{warnbox}[版本会变]
本章的命令以 0.2.2-rc.5 为准。若安装到的版本更新，参数名与配置文件路径都可能不同，以 \texttt{dsh -{}-help} 与实际生成的配置为准，不要照抄教程。
\end{warnbox}

\section{安装}

三种方式，按可复现程度排列。

\begin{shellsession}
$ npx @deepseek-ai/dsh web          # 不安装，直接运行
$ npm install -g @deepseek-ai/dsh   # 全局安装
$ nix shell nixpkgs#nodejs_24 --command \
    npx @deepseek-ai/dsh --version
\end{shellsession}

在 NixOS 上，全局安装到系统会污染 profile，推荐写进配置：

\begin{nixcode}
environment.systemPackages = with pkgs; [
  nodejs_24
];

# 或用临时环境，需要时进入
# $ nix shell nixpkgs#nodejs_24
\end{nixcode}

\begin{shellsession}
$ nix shell nixpkgs#nodejs_24
$ npm install -g @deepseek-ai/dsh
$ dsh --help
\end{shellsession}

\section{第一次启动}

\begin{shellsession}
$ dsh web
$ dsh web --no-open       # 不自动打开浏览器
\end{shellsession}

Web 界面默认监听 \texttt{127.0.0.1:3080}，在浏览器打开即可。远程服务器上使用时，不要直接对外开放该端口，改用 SSH 端口转发：

\begin{shellsession}
$ ssh -L 3080:127.0.0.1:3080 srv1
\end{shellsession}

本地浏览器访问 \texttt{localhost:3080} 即等于访问服务器上的界面，端口不暴露在公网。

\section{模型与凭据}

模型接入在设置里配置。两件事必须分清：Agent 程序本身与它调用的模型提供方是两回事，安装 Agent 不等于获得模型访问权。

\begin{booktable}{\textwidth}{J{0.815} J{1.185}}
\toprule
概念 & 内容 \\
\midrule
\endhead
提供方 & 模型的来源，例如官方 API 或第三方兼容接口 \\
凭据 & API 密钥，决定用量与计费 \\
模型 & 具体的模型名，不同模型能力与价格不同 \\
\bottomrule
\end{booktable}

密钥不要写进任何会被提交的文件。放在环境变量或 \texttt{/run/secrets} 下，由启动脚本注入。第三方兼容接口需要填写地址与协议，协议填错的表现是请求 404 或 400，而不是鉴权失败。

\section{权限与沙箱}

工具调用能执行命令，因此默认应当收窄权限。首次运行会在交互界面询问是否允许某类操作，选择时按最小必要判断：读文件可以允许，删除与网络上传应当逐次确认。

\begin{shellsession}
$ dsh --help | grep -i -E "sandbox|approval|permission"
\end{shellsession}

若版本提供了仅规划不执行的模式，写复杂任务前先用它跑一遍，看清计划再放开执行。这一步比事后回滚便宜。

\section{插件与扩展}

Harness 的架构是“一切皆插件”。常见扩展方式有两种：安装官方或社区插件，或通过 MCP 接入外部工具服务。安装插件前确认三件事：来源仓库、是否需要额外凭据、能否卸载。插件与 Agent 同权限运行，因此来源不明的插件等同于来源不明的程序。

\begin{shellsession}
$ dsh plugin list
$ dsh plugin --help
\end{shellsession}

\section{一次完整任务}

目标：在仓库里找出所有硬编码的时区字符串，改成从环境变量读取。

\begin{enumerate}
  \item 在项目目录启动：\texttt{cd ~/proj \&\& dsh web}。
  \item 描述目标与约束：范围限于 \texttt{src/}，不改测试，改完跑一次测试。
  \item 观察工具调用：它应当先搜索、再读文件、再改。
  \item 检查改动：用 \texttt{git diff} 逐段看，确认没有无关改动。
  \item 跑验证：测试、构建、静态检查，任何一项失败都退回第一步补充约束。
\end{enumerate}

\begin{shellsession}
$ git status --short
$ git diff
$ npm test
\end{shellsession}

三条命令是验收的标准动作。Agent 说“已完成”不算完成，差异与测试结果才算。

\section{练习}

\begin{enumerate}
  \item 用 \texttt{npx} 运行一次并记录版本号，与本章标注的版本比较。
  \item 用 \texttt{-{}-no-open} 启动，通过 SSH 转发访问界面，确认端口未对外暴露。
  \item 配置一个模型提供方，发一条最短的请求，记录耗时。
  \item 用只读权限完成一次代码检索任务，确认它没有修改任何文件。
  \item 完成上面的一次完整任务，用 \texttt{git diff} 与测试结果验收。
\end{enumerate}

\begin{checkpoint}
本章结束时，应能：
\begin{itemize}
  \item 说出当前版本、运行时要求与许可；
  \item 用三种方式安装或运行，知道 NixOS 上的推荐做法；
  \item 通过 SSH 转发安全地使用 Web 界面；
  \item 区分 Agent、提供方、凭据与模型四个概念；
  \item 按五步流程完成一次任务并用差异与测试验收。
\end{itemize}
\end{checkpoint}
